\chapter{عبارت‌های جبری و چندجمله‌ای‌ها}
\thispagestyle{empty}

\begin{quote}
	{\lotus
		جبر سخاوتنمد است؛ اغلب بیش از آنچه از وی تقاضا می‌کنیم به ما می‌بخشد. \quad - دالامبر%
		\LTRfootnote{ Jean Le Rond d'Alembert}
	}
\end{quote}

به سه مسئله‌ای که در ادامه می‌آید، توجه کنید. شباهت و تفاوت آن‌ها در چیست؟ \\
۱. ۱۲ سیب را بین دو نفر طوری تقسیم کنید که یکی از آن‌ها دو برابر دیگری داشته باشد. \\
۲. عدد ۱۲ را به نسبت ۱ به ۲ تقسیم کنید. \\
۳. عدد $a$ را به نسبت $n$ به $m$ تقسیم کنید.

\section{عبارت‌های جبری}
در ریاضیات برای نشان دادن یک متغیر از نمادهایی که غالباً حروف لاتین و انگلیسی است، استفاده می‌شود و به عبارت ایجادشده یک عبارت جبری گفته می‌شود.
\begin{example}
	در زیر چند عبارت جبری نوشته شده است.
	\begin{gather*}
		(a+b)^3 ; \\
		2x + \sqrt{x-1} \\
		\frac{2x^2-x-7}{3x^2-5}
	\end{gather*}
\end{example}

\subsection{عبارت‌های جبری گویا و گنگ}
\begin{example}
	عبارت‌های
	$(a+b)^3$، $5x^2-3x+\sqrt{2}$ و $\frac{2x-1}{3x-5}$
	عبارت‌های جبری گویا هستند.
\end{example}

\begin{example}
	عبارت‌های
	$\sqrt{5xy}$، $2x+\sqrt{x-1}$ و $\frac{5x\sqrt{x+1}}{x+2}$
	عبارت‌های جبری گنگ هستند.
\end{example}

\section{جمله}
عبارت‌هایی مانند
$2x^2$، $-3y$، $4xy^2$ و $ax^n$ ($n$ عددی صحیح و مثبت یا صفر)،
هر کدام یک جمله را معرفی می‌کنند و در جبر به آن‌ها یک‌جمله‌ای یا جمله می‌گویند. جمله یعنی عبارتی که در آن برای حرف‌ها تنها نماد ضرب و توان طبیعی داشته باشد. حاصل‌ضرب عامل‌های عددی و معلوم را ضریب عددی یک‌جمله‌ای و مجموع نماهای متغیرهای آن را درجه‌ی یک‌جمله‌ای می‌گویند. ضریب عددی می‌تواند هر عدد حقیقی باشد. 

\begin{example}
	یک‌جمله‌ای $5x^2y^3z$ را در نظر بگیرید. ضریب عددی این یک‌جمله‌ای ۵ است. این یک جمله‌ای نسبت به $x$ از درجه‌ی ۲، نسبت به $y$ از درجه‌ی ۳، نسبت به $z$ از درجه‌ی یک و نسبت به سایر حروف از درجه‌ی صفر است. همچنین درجه‌ی یک‌جمله‌ای نسبت به $x$ و $y$ و $z$ برابر $2+3+1=6$ است.
\end{example}

\begin{example}
	عبارت‌های
	$\frac{4xy}{z}$، $5x^2y^2+1$ و $\sqrt{2x}$
	یک‌جمله‌ای نیستند.
\end{example}

\subsection{جمله‌های متشابه}
دو جمله اگر تنها در ضریب عددی خود اختلاف داشته باشند، متشابه نامیده می‌شوند.
\begin{example}
	جمله‌های
	$\frac{1}{2}x$، $3x$ و $x\sqrt{5}$
	متشابه هستند.
\end{example}

\begin{example}
	دو جمله‌ی
	$4a^3b$
	و $5a^3b^2$
	نسبت به دو متغیر $a$ و $b$ متشابه نیستند؛ ولی نسبت به متغیر $a$ متشابه‌اند.
\end{example}

\subsection{جمع و تفریق یک‌جمله‌ای جبری}
دو یا چند یک‌جمله‌ای را در صورتی می‌توان جمع و تفریق کرد که جمله‌ها متشابه باشند. برای این منظور باید ضریب‌های عددی را با هم جمع یا تفریق کرد.
\begin{example}
	\begin{gather*}
		7ab + 3ab = 10ab ; \\
		\frac{2}{5}x^3 + x^3 = \frac{7}{5}x^3 ; \\
		x^2y - ax^2y + 3ax^2y = (2a+1)x^2y
	\end{gather*}
\end{example}

\subsection{ضرب و تقسیم یک‌جمله‌ای جبری}
برای انجام عمل ضرب یا تقسیم در یک‌جمله‌ای‌ها ضریب‌های عددی را در هم و متغیرهای هم‌نام را در یکدیگر ضرب یا بر هم تقسیم می‌کنیم.
\begin{example}
	\begin{gather*}
		5x^2y^3 \times 2x^3y^2z = 10x^5y^5z ; \\
		\frac{15x^2y}{3xy^2} = \frac{5x}{y} ; \\
		2a^2b \times 3a^3b^4 + 4a \times a^4b^5 = 6a^5b^5 + 4a^5b^5 = 10a^5b^5
	\end{gather*}
\end{example}

\begin{example}
	اگر
	$A = 6x^2y^3z^4$
	و
	$B = 4x^4y^3z^2$
	و
	$C = 12x^2y^2z^2$،
	حاصل عبارت
	$\frac{A \times 2B}{C}$
	را به دست آورید. \\
	\begin{equation*}
		\frac{A \times 2B}{C} = \frac{6x^2y^3z^4 \times 2 \times 4x^4y^3z^2}{12x^2y^2z^2} = \frac{48x^6y^6z^6}{12x^2y^2z^2} = 4x^4y^4z^4
	\end{equation*}
\end{example}

\subsection{بخش‌پذیری یک‌جمله‌ای بر یک‌جمله‌ای}
یک‌جمله‌ای $A$ بر یک‌جمله‌ای $B$ بخش‌پذیر است، در صورتی که حاصل تقسیم $A$ بر $B$ خود یک‌جمله‌ای باشد.

\section{چندجمله‌ای}

\section{اتحادها}

\begin{equation}
	(a+b)^2 = a^2+2ab+b^2
\end{equation}

\begin{equation}
	(a-b)^2 = a^2-2ab+b^2
\end{equation}

\begin{equation}
	a^2-b^2 = (a+b)(a-b)
\end{equation}

\begin{equation}
	(x+a)(x+b) = x^2+(a+b)x+ab
\end{equation}

\begin{equation}
	(a+b)^3 = a^3+3a^2b+3ab^2+b^3
\end{equation}

\begin{equation}
	(a-b)^3 = a^3-3a^2b+3ab^2-b^3
\end{equation}

\begin{equation}
	a^3+b^3 = (a+b)(a^2-ab+b^2)
\end{equation}

\begin{equation}
	a^3-b^3 = (a-b)(a^2+ab+b^2)
\end{equation}

\begin{equation}
	(a+b+c)^2 = a^2+b^2+c^2+2ab+2bc+2ca
\end{equation}

\begin{equation}
	(a+b+c)^3 = a^3+b^3+c^3+3(a+b)(b+c)(c+a)
\end{equation}

\begin{equation}
	a^3+b^3+c^3 -3abc = (a+b+c)(a^2+b^2+c^2-ab-bc-ca)
\end{equation}

\begin{equation}
	a^3+b^3+c^3 -3abc = \frac{1}{2}(a+b+c)\left((a-b)^2+(b-c)^2+(c-a)^2\right)
\end{equation}

\begin{equation}
	(a^2+b^2)(c^2+d^2) = (ac+bd)^2 + (ad-bc)^2
\end{equation}

\begin{equation}
	(a^2+b^2)(c^2+d^2) = (ac-bd)^2 + (ad+bc)^2
\end{equation}

\begin{example}
	با فرض
	$x+\dfrac{1}{x}=a \geq 2$
	حاصل عبارت
	$x^3 -\dfrac{1}{x^3}$
	را برحسب $a$ به دست آورید. \\
	\textbf{پاسخ:}
	\[
	x+\frac{1}{x}=a \Rightarrow \left(x+\frac{1}{x}\right)^2=a^2 \Rightarrow x^2+2+\frac{1}{x^2}=a^2 \Rightarrow x^2+\frac{1}{x^2}=a^2-2
	\]
	\[
	\left(x-\frac{1}{x}\right)^2=x^2+\frac{1}{x^2}-2=\left(a^2-2\right)-2=a^2-4 \Rightarrow x-\frac{1}{x}=\pm \sqrt{a^2-4}
	\]
	\[
	x^3-\frac{1}{x^3}=\left(x-\frac{1}{x}\right)\left(x^2+1+\frac{1}{x^2}\right) = \left(\pm \sqrt{a^2-4}\right)\left(a^2-1\right)
	\]
\end{example}
